% ID: SEL-D0059-OA-E03
% ESTADO: cotejada. Segunda pasada visual y borrador compilado revisados.
% Fuente: Selectividad/sel_2011_matematicas_aplicadas/222-EXAMEN 2.pdf
% SHA-256: 998f20402054ca6d0454da83ffc8e1ff6951a576b48836a32e81e524c5985214
% Página: 1; opción A; ejercicio original 3.

Un jugador lanza a la vez un dado y una moneda.
\begin{enumerate}[label=\alph*)]
\item \textbf{(1 punto)} Construya el espacio muestral de este experimento aleatorio.
\item \textbf{(1 punto)} Determine la probabilidad del suceso $A$: ``El jugador obtiene un número par en el dado y cruz en la moneda''.
\item \textbf{(0.5 puntos)} Si sabemos que en la moneda ha salido cara, \textquestiondown{}cuál es la probabilidad de que en el dado haya salido más de 3 puntos?
\end{enumerate}
